\section{Reasoning Agent 的理论意义}

\subsection{推理作为无限维的内部动作}

传统 Agent（无论 symbolic 还是 RL）的动作空间由环境固定定义（如上下左右），而 Reasoning Agent 引入了一种全新的动作类型——\textbf{推理}：

\begin{itemize}
    \item 推理可以是任意长度的自然语言，动作空间是\textbf{无限的}
    \item 推理不改变外部环境，只改变 Agent 自身的上下文（内部状态）
    \item 推理影响后续的行动决策
\end{itemize}

\subsection{三种 Agent 范式的对比}

\begin{table}[H]
\centering
\begin{tabular}{lccc}
\toprule
& \textbf{Symbolic AI Agent} & \textbf{Deep RL Agent} & \textbf{Reasoning Agent} \\
\midrule
内部表示 & 符号状态 & 神经嵌入（向量） & 自然语言 \\
构建方式 & 手写规则 & 训练数百万步 & Prompting / Fine-tuning \\
领域迁移 & 极差 & 差 & 好 \\
表示维度 & 有限符号集 & 固定维度向量 & 无限语言空间 \\
\bottomrule
\end{tabular}
\caption{三种 Agent 范式对比}
\end{table}

自然语言作为中间表示的优势：
\begin{itemize}
    \item 利用 LLM 预训练的丰富知识，无需大量额外训练
    \item 表达能力无限——可以思考一个词，也可以思考一百万个 token
    \item 天然支持 inference-time scaling
\end{itemize}

\subsection{本章小结}

Reasoning Agent 的本质创新在于用自然语言作为 Agent 的内部表示和推理媒介，这赋予了 Agent 无限维的内部动作空间和预训练知识的零样本迁移能力。

